\thispagestyle{empty}
\section*{\fontschap Resumen \footnote{¿Cómo citar este libro? / How to cite this book?
		Zambrano Caviedes, J. y Baquero Torres, E. (2025). Problemas resueltos de integración múltiple con aplicaciones en Wolfram Mathematica. Universidad Distrital Francisco José de Caldas-Editorial UD. DOI: XXXX}}

Presentamos en estas notas de clase, algunos problemas 
rutinarios sobre integración para un curso de Cálculo Multivariado. En ellas se trabaja, básicamente, la integración de una función sobre un objeto que puede ser una curva, superficie o sólido, con el fin de parametrizar la región de integración.
Esta idea es diferenciadora con respecto a la mayoría de textos de cálculo multivariado. Cada problema ofrece un gráfico asociado en diversas vistas y el cálculo que se requiere. Además se incluyen las instrucciones utilizadas en el programa  \emph{Wolfram Mathematica}. El cambio de sistema de coordenadas para evaluar integrales triples, se hace tradicionalmente utilizando coordenadas esféricas y cilindricas; no obstante, se incluyen otros sistemas de coordenadas en los anexos y una breve referencia en los ejercicios propuestos para motivar al lector al estudio de estos tópicos.

\fontschap Palabras clave:\normalfont integración múltiple, parametrización, campos escalares, campos vectoriales, Wolfram Mathematica.
\vspace{-0.7em} % Espacio antes de la línea

\noindent
\hspace*{-0.1cm} % Ajusta este valor para mover la línea hacia la derecha
\rule{2cm}{0.5pt} % Línea horizontal

\vspace{-0.5em} % Espacio después de la línea



\section*{\fontschap Abstract}

We present in these class notes some problems 
routines on integration for a multivariate calculus course. In them, we work, basically, integrating a function on an object that can be a curve, surface or solid, parameterizing the region of integration.
This idea is different from most multivariate calculus texts. Each problem offers an associated graph in various views and the calculation that is required, the instructions used in the {\tt Wolfram Mathematica} program are also included. Changing the coordinate system to evaluate triple integrals is traditionally done using spherical and cylindrical coordinates. However, other coordinate systems are included in the annexes and a brief reference in the proposed exercises to motivate the reader to study these topics.

\fontschap Keywords:\normalfont multiple integration, parameterization, scalar fields, vector fields, Wolfram Mathematica.

\begin{spacing}{1.4}
	\begin{tabular}{p{0.26\textwidth}!{\vrule width 0.9pt}p{0.69\textwidth}}
		& \hspace{2mm}{\medseries \fontsize{11pt}{9pt}\selectfont{¿Como citar este libro?~/ How to cite this book?}} \\[3.5pt]
		& \hangindent=2mm \hangafter=0
		\fontsize{9pt}{9pt}\selectfont{Salamanca Bernal, J. A., Rodríguez Patarroyo, D. J. y Pantoja Benavides, J. F. (2025). \emph{Columna total de ozono estratosférico y manchas solares. Una relación físico-estadística mundial desde la minería de datos geoespaciales NASA}. Universidad Distrital Francisco José de Caldas-Editorial UD. https://doi.org/10.69740/UD.9789587877724.9789587877748}\\
	\end{tabular}
\end{spacing}